\section{Methodology}
\label{sec:methodology}

\subsection{The Solution Space and Why Seeds Matter}
\label{sec:methodology:space}

Graph minimum $k$-bisection is NP-hard. METIS finds a local optimum via
multilevel coarsening and refinement; the random seed controls the initial
partitioning from which refinement proceeds. Critically, \emph{every seed
produces a distinct final partition} in the high-dimensional partition space
--- we verified this empirically for Pennsylvania across 100 seeds (100 unique
assignment hashes, no two identical). The solution space is therefore dense and
continuous, not a small collection of canonical boundaries.

The partisan outcome of a plan is quantized by the integer seat constraint:
with $k$ districts and statewide Dem share $v$, the achievable proportionality
gaps are $\{100(j/k - v) : j = 0, 1, \ldots, k\}$. This quantization explains
why gap values appear to cluster --- different boundaries can produce the same
integer seat count --- while the underlying plan space is effectively infinite.

\subsection{The Minimum-Edge-Cut Protocol}
\label{sec:methodology:protocol}

Rather than report results for a single arbitrary seed, we propose the
\emph{minimum-edge-cut (MEC) protocol}: run $N$ seeds, compute the total
boundary length (edge-cut) of each resulting partition, and publish the
plan with the smallest edge-cut. Formally, for state $s$ and seed set
$\mathcal{S} = \{1, \ldots, N\}$:

\[
  \pi^*_s = \arg\min_{\pi \in \{\pi_i : i \in \mathcal{S}\}} \mathrm{EC}(\pi)
\]

where $\mathrm{EC}(\pi) = \sum_{(u,v) \in E,\, \pi(u) \neq \pi(v)} w_{uv}$
is the total TIGER boundary length (in metres) of cut edges. This protocol is:
\begin{itemize}
  \item \textbf{Deterministic given $\mathcal{S}$}: publishing the full seed
        set allows anyone to reproduce $\pi^*_s$ exactly.
  \item \textbf{Monotone}: adding seeds can only decrease or maintain
        $\mathrm{EC}(\pi^*_s)$.
  \item \textbf{Parti\-san-blind}: no partisan data enters the selection
        criterion; the choice is made entirely on geometric grounds.
  \item \textbf{Convergence-testable}: we can measure how quickly the minimum
        stabilises across seeds, and declare ``enough seeds'' when the
        probability of improvement drops below a calibrated threshold.
\end{itemize}

\subsection{Convergence Analysis}
\label{sec:methodology:convergence}

Let $M_n = \min_{i \leq n} \mathrm{EC}(\pi_i)$ be the running minimum after
$n$ seeds. We model the edge-cut values as i.i.d.\ draws from an unknown
distribution $F$ on $[\mathrm{EC}_{\min}, \infty)$ and estimate $F$
non-parametrically from the observed sample. The probability that seed $n+1$
sets a new minimum is $P(X < M_n) = F(M_n^-)$, which decreases as $n$
increases (since $M_n$ is monotone decreasing).

We declare convergence when the empirical probability of improvement over the
next 100 seeds falls below 1\% --- equivalently, when the observed gap between
the $n$th-best and $(n+1)$th-best is consistent with the tail of a fitted
extreme-value distribution. We validate convergence separately for each state.

\emph{Empirical calibration (Pennsylvania, 100 seeds):} The running minimum
improved 7 times at seeds $\{1,2,3,5,12,18,34\}$, then stabilised for seeds
35--100. The improvement rate at $n=100$ is $\hat\lambda \approx 0.02$
improvements per seed, suggesting roughly 5,000 seeds would be needed to
achieve a 90\% confidence interval of $\pm$1\,km on the true minimum for PA.
For states with fewer districts (e.g., Vermont with 1 district), the true
minimum is trivially achieved at every seed; for large states (Texas, 38
districts), convergence will be slower.

\subsection{Exhaustive Seed Sweep Design}
\label{sec:methodology:sweep}

Based on the Pennsylvania calibration, we run the following sweep:

\begin{itemize}
  \item \textbf{Phase 1 (deep dive)}: Vermont ($k=1$), Pennsylvania ($k=17$),
        Texas ($k=38$) $\times$ 5{,}000 seeds each. Goal: characterise full
        convergence curves and establish the partisan-outcome distribution over
        the near-MEC solution space.
  \item \textbf{Phase 2 (50-state)}: All 50 states $\times$ 1{,}000 seeds
        each. Goal: identify which states have seed-sensitive partisan outcomes
        and what fraction of the near-MEC space each outcome class occupies.
  \item \textbf{Phase 3 (convergence validation)}: For each state in Phase 2,
        fit an extreme-value (GEV) model to the observed minimum sequence and
        report the 95\% CI on the estimated true minimum. States where the CI
        spans more than one seat boundary are flagged as ``not yet converged''
        and receive additional seeds.
\end{itemize}

Compute estimate: $50 \times 1{,}000 = 50{,}000$ state-level runs at an
average of $\sim$30\,s each $\approx$ 417 CPU-hours. Parallelised across 8
cores: $\sim$52 wall-clock hours.

\subsection{CompactBisect: A Huntington-Hill Analogue for Geometric Redistricting}
\label{sec:methodology:compactbisect}

\subsubsection{Why Compactness Supersedes Edge-Cut as the Selection Criterion}

Minimum edge-cut with boundary-length weights minimises total district
perimeter contribution from \emph{internal} boundaries --- the shared edges
cut between districts. But the Polsby-Popper score of a district is
$\mathrm{PP}_i = 4\pi A_i / P_i^2$, where $P_i$ includes \emph{both} the
internal boundaries (controlled by edge-cut) and the \emph{external} state
boundary: coastlines, rivers, and state lines that are geographically fixed
regardless of how the partition is drawn.

For coastal and complex-border states --- Maryland, Florida, California,
Louisiana --- a district along the water's edge will have a long external
perimeter that no algorithm can reduce. Edge-cut minimisation is blind to
this; Polsby-Popper is not. Compactness directly measures the geometric
quality we want: how close to a circle is each district, accounting for
its full boundary. Edge-cut is a computationally efficient heuristic for
generating candidate bisections; compactness is the correct selection criterion.

\subsubsection{The Geometric Mean as the Fairness Criterion}

Given a candidate bisection $(L, R)$ of subgraph $G$, three natural scores
summarise the compactness of both halves:

\begin{itemize}
  \item \textbf{Minimum}: $\min(\mathrm{PP}(L), \mathrm{PP}(R))$ ---
        maximise the worst half (Rawlsian).
  \item \textbf{Arithmetic mean}: $(\mathrm{PP}(L) + \mathrm{PP}(R))/2$ ---
        maximise total compactness.
  \item \textbf{Geometric mean}: $\sqrt{\mathrm{PP}(L) \cdot \mathrm{PP}(R)}$
        --- maximise \emph{balance} between the two halves.
\end{itemize}

The geometric mean is strictly superior for the following reason. For a fixed
arithmetic mean $\mu = (\mathrm{PP}(L) + \mathrm{PP}(R))/2$, the geometric
mean $\sqrt{\mathrm{PP}(L)\cdot\mathrm{PP}(R)} \leq \mu$ with equality iff
$\mathrm{PP}(L) = \mathrm{PP}(R)$. Maximising the geometric mean therefore
simultaneously maximises compactness \emph{and} enforces equality between the
two halves. It penalises imbalance in proportion to its magnitude: a plan
with $\mathrm{PP}(L) = 0.1$, $\mathrm{PP}(R) = 0.9$ scores 0.30 (geometric)
vs.\ 0.50 (arithmetic). A plan with $\mathrm{PP}(L) = \mathrm{PP}(R) = 0.5$
scores 0.50 under both --- the same, as it should be.

\paragraph{The Huntington-Hill analogy.}
This criterion is structurally identical to the Huntington-Hill equal-proportions
method for congressional apportionment (2~U.S.C.~\S~2a). Huntington-Hill
assigns the next seat to the state for which
$\sqrt{n_i(n_i+1)}$ is minimised, where $n_i$ is the current seat count. The
geometric mean $\sqrt{n_i(n_i+1)}$ is the breakpoint because it enforces
proportional equality in representation: it treats a move from $n$ to $n+1$
seats proportionally the same regardless of $n$. Congress adopted this
criterion in 1941 precisely because it is the uniquely neutral geometric-mean
fairness standard for apportionment.

CompactBisect applies the identical reasoning to the bisection problem.
Rather than proportional equality in seats, we seek proportional equality in
compactness: $\sqrt{\mathrm{PP}(L)\cdot\mathrm{PP}(R)}$ is the
\emph{geometric-mean compactness} of the two halves, and maximising it ensures
that neither half is penalised for being proportionally less compact than the
other. The federal government already accepts geometric-mean fairness as the
uniquely correct apportionment standard; CompactBisect applies the same
principle one level down, to the line-drawing problem itself.

\subsubsection{Algorithm Definition}

\begin{definition}[\textsc{CompactBisect}]
\label{def:compactbisect}
Let $G = (V, E, w_e)$ be the state census-tract graph with edge weights equal
to TIGER shared boundary lengths (metres), per-vertex areas $a_v$ (m$^2$), and
per-vertex external perimeters $p_v^{\mathrm{ext}}$ (metres, the non-shared
boundary of each tract). For a subset $S \subseteq V$, define:
\[
  A(S) = \sum_{v \in S} a_v, \qquad
  P(S) = \underbrace{\sum_{(u,v)\in E,\, u\in S,\, v\notin S} w_e^{(uv)}}_{\text{internal boundary}}
        + \underbrace{\sum_{v\in S} p_v^{\mathrm{ext}}}_{\text{external boundary}},
  \qquad
  \mathrm{PP}(S) = \frac{4\pi A(S)}{P(S)^2}.
\]

\textsc{CompactBisect}$(G, k, N, \varepsilon)$, with $N$ seeds and
$\varepsilon$-tolerance ($\varepsilon = 0.05$ unless stated):
\begin{enumerate}
  \item If $k = 1$, return $\{V\}$ as a single district.
  \item For $i = 1, \ldots, N$: run METIS on $G$ with seed $i$ to obtain
        bisection $(L_i, R_i)$ with target sizes
        $(\lfloor k/2\rfloor, \lceil k/2\rceil)$.
        Record edge-cut $\mathrm{EC}_i = \sum_{u\in L_i,\, v\in R_i} w_e^{(uv)}$.
  \item Let $\mathrm{EC}^* = \min_i \mathrm{EC}_i$.
        (METIS serves purely as the candidate generator; edge-cut is the
        search heuristic, not the selection criterion.)
  \item Select:
        \[
          i^* = \arg\max_{i:\, \mathrm{EC}_i \leq (1+\varepsilon)\mathrm{EC}^*}
                \sqrt{\mathrm{PP}(L_i)\cdot\mathrm{PP}(R_i)}.
        \]
  \item Recurse: \textsc{CompactBisect}$(G[L_{i^*}],\,\lfloor k/2\rfloor,\,N,\,\varepsilon)$
        $\cup$ \textsc{CompactBisect}$(G[R_{i^*}],\,\lceil k/2\rceil,\,N,\,\varepsilon)$.
\end{enumerate}
\end{definition}

\subsubsection{Properties}

\paragraph{Level-by-level justifiability.}
Any plan produced by \textsc{CompactBisect} can be justified at every
intermediate step: at each bisection, the algorithm chose the candidate
maximising geometric-mean compactness among near-minimum-cut splits.
No intermediate decision can be characterised as partisan because the
selection criterion is purely geometric and admits no partisan input.
This is directly analogous to the statutory justification for
Huntington-Hill apportionment: the algorithm is specified in law, the
inputs are public census figures, and the output is deterministic.

\paragraph{Compactness equity across parties.}
A plan may be compact on average but with systematically less compact
districts for one party's coalition (a residual signature of
packing-and-cracking). \textsc{CompactBisect}'s geometric-mean criterion
enforces equal compactness at each bisection level, which propagates to
the final districts. We test the residual correlation empirically using
a Spearman rank correlation between final-district Polsby-Popper scores
and district Dem two-party vote shares; a neutral plan should show
$|\rho| < 0.1$.

\paragraph{Computational cost.}
With $k$ districts, $\ell = \lceil\log_2 k\rceil$ recursion levels,
and $N$ seeds per level, \textsc{CompactBisect} performs $N\ell$ METIS
calls total. For Pennsylvania ($k=17$, $\ell=5$, $N=100$): 500 calls
$\approx$ 5~min. For Texas ($k=38$, $\ell=6$, $N=100$): 600 calls
$\approx$ 10~min. Cost scales as $O(N\log k)$, not $O(Nk)$.

\paragraph{Relationship to standard edge-weighted bisection.}
When $N=1$ and $\varepsilon=0$, \textsc{CompactBisect} degenerates to
a single METIS run with one fixed seed --- the current default algorithm.
As $N\to\infty$, the candidate pool densely covers the near-minimum-cut
solution space and the geometric-mean compactness criterion selects the
globally optimal plan within it. The parameter $N$ controls the
exploration-exploitation tradeoff; we calibrate $N$ via convergence
analysis (\S\ref{sec:methodology:convergence}).

\subsection{The Proportionality Metric}
\label{sec:methodology:metric}

The primary outcome variable is the \emph{proportionality gap}:
\[
  \Delta_s = 100 \cdot \left(\frac{d_s}{k_s} - v_s\right)
\]
where $d_s$ is the number of Dem-leaning districts (Dem two-party share
$\geq 0.5$), $k_s$ is the total district count, and $v_s$ is the statewide
Dem two-party vote share from the 2020 presidential election. Positive
$\Delta_s$ indicates Democratic over-representation; negative indicates
Republican over-representation. We computed $\Delta_s$ for all 50 states at a
single reference seed as baseline (Task 1 results).

The key question is whether the MEC protocol --- taking the minimum-edge-cut
plan across $N$ seeds --- systematically shifts $\Delta_s$ in a partisan
direction relative to the single-seed baseline and relative to the full
seed-distribution mean.
